\chapter{Metodología de Evaluación}
\label{ch:metodologia}

En este capítulo se establece el diseño de la investigación aplicada, los criterios de delimitación del alcance, el catálogo de herramientas y el protocolo de auditoría técnica no destructiva \citep{nist2008sp800115}.

\section{Tipo y Nivel de Investigación}
El presente trabajo corresponde a una \textbf{investigación aplicada}, orientada a la resolución de problemas de seguridad concretos en una infraestructura real. Posee un nivel \textbf{descriptivo y analítico}, integrando métodos cualitativos (para el modelado y categorización STRIDE) y cuantitativos (en la evaluación del impacto y probabilidad de riesgos con DREAD y CVSS v3.1) \citep{first2019cvss,meier2003improving}.

\section{Área de Estudio y Perímetro Institucional}
La investigación se circunscribe a la infraestructura de servicios web públicos e internos de la \textbf{Universidad Nacional del Altiplano de Puno} (\texttt{unap.edu.pe}):

\begin{itemize}
    \item \textbf{Portal Académico SIA (\texttt{sia.unap.edu.pe}):} Gestión de actas, notas oficiales y matrículas.
    \item \textbf{Campus Virtual / LMS (\texttt{vcampus.unap.edu.pe}):} Plataforma Moodle para docencia y evaluaciones.
    \item \textbf{Portal de Admisión (\texttt{admision.unap.edu.pe}):} Inscripciones de postulantes y pagos.
    \item \textbf{API de Servicios Institucionales (\texttt{api.unap.edu.pe}):} Interconexión REST con la app móvil y bases de datos centrales.
\end{itemize}

\section{Catálogo de Herramientas y Propósitos de Evaluación}
Para garantizar la reproducibilidad y rigor técnico del estudio, se seleccionó un conjunto de herramientas estándar de auditoría y análisis de seguridad, clasificadas según su propósito operativo:

\begin{center}
\small
\begin{longtable}{>{\bfseries}p{2.8cm} p{4.2cm} p{4.5cm} p{3.2cm}}
\toprule
\rowcolor{outermarginbgcolor}
Herramienta & Propósito Técnico Principal & Técnica / Estándar Asociado & Fase de Aplicación \\
\midrule
\endfirsthead
\toprule
\rowcolor{outermarginbgcolor}
Herramienta & Propósito Técnico Principal & Técnica / Estándar Asociado & Fase de Aplicación \\
\midrule
\endhead
\bottomrule
\endfoot
whois & Consulta de registros de asignación IP y contactos del dominio institucional. & OSINT / T1596 (Search Open Technical Databases) & Reconocimiento (Kill Chain F1) \\
\midrule
dig / host & Inspección de registros autoritativos DNS (A, NS, MX, TXT) y transferencias de zona. & Reconocimiento DNS / T1590 (Gather Network Info) & Reconocimiento (Kill Chain F1) \\
\midrule
nmap & Detección de puertos abiertos, servicios activos y fingerprinting de versiones. & Active Scanning / T1595.001 (Port Scans) & Reconocimiento (Kill Chain F1) \\
\midrule
ssh-audit & Auditoría de algoritmos de cifrado, intercambio de claves y MAC en servicio SSH. & T1190 (Exploit Public-Facing Application) & Reconocimiento / Análisis Criptográfico \\
\midrule
curl & Inspección de cabeceras HTTP de respuesta, directivas de seguridad y TLS. & Fingerprinting Web / OWASP WSTG-INFO-07 & Info Disclosure (STRIDE / F1) \\
\midrule
sublist3r & Enumeración pasiva de subdominios a través de motores de búsqueda y certificados. & OSINT Subdomain Enumeration & Reconocimiento (Kill Chain F1) \\
\midrule
SQLMap & Simulación de inyección SQL automatizada y validación de sanitización de entradas. & T1190 (SQL Injection) / Tampering & Explotación (Kill Chain F4 / STRIDE) \\
\midrule
Burp Suite & Intercepción de tráfico HTTP, manipulación de parámetros y análisis de tokens. & Broken Access Control / IDOR / Spoofing & Explotación / Elevación Privilegios \\
\midrule
slowhttptest & Simulación de agotamiento de sockets HTTP Layer 7 (ataque Slowloris). & T1499.003 (Endpoint DoS) / Denial of Service & Explotación (Kill Chain F4 / STRIDE) \\
\midrule
Wazuh SIEM & Agregación de logs, correlación de eventos y detección de firmas de intrusión. & Blue Team SIEM Detection / NIST CSF & Propuesta de Mitigación y Monitoreo \\
\bottomrule
\end{longtable}
\end{center}

\section{Procedimiento Metodológico en Tres Fases}
El procedimiento se divide en tres fases rigurosamente secuenciales:
\begin{enumerate}
    \item \textbf{Fase 1: Reconocimiento Pasivo y Recolección de Datos:} Identificación de activos expuestos, resolución de infraestructura y análisis de configuraciones perimetrales mediante técnicas OSINT no invasivas.
    \item \textbf{Fase 2: Modelado Sistemático STRIDE por Activo:} Desglose arquitectónico de cada servicio web y mapeo contra las seis categorías de amenazas de Microsoft.
    \item \textbf{Fase 3: Reconstrucción Operativa Cyber Kill Chain y Detección:} Simulación controlada de los 3 escenarios de mayor impacto a través de las 7 etapas de intrusión, formulando contramedidas específicas de ruptura de cadena.
\end{enumerate}
